% ============================================================
% 分片：s05c-topology  对应原文 PDF 第 66--74 页（印刷页码 66--74），另收第 65 页末
%       跨页半句与第 75 页顶部 3 行（§5.5 末句，见下方边界说明）
% 负责章节：5.5 Change Of Topology
% 译者：t06-s33   审校状态：初稿
% 审校：t04-s3a 已完成逐式核对（修正 3 处）
% ============================================================
% 边界说明：
%   1) §5.4 末段的句子 "First of all, a suitable quotient of the variety X by a finite
%      group gives one of the first known and simplest three | generation models,
%      constructed originally by Schimmrigk [42]." 跨第 65/66 页。s05b 止于其前一句
%      "This model is of interest for two reasons."，故本片自第 65 页末那半句起笔，
%      再接第 66 页首 "generation models, constructed ..."。该段中的公式 (5.36) 属 §5.4。
%   2) 第 74 页末句 "…have a pole at t_a = 0 as" 跨至第 75 页顶部 3 行，之后即
%      "6. EXTENSION TO (0,2) MODELS" 标题。为保持 §5.5 完整，第 75 页顶部 3 行
%      （§5.5 末句）收入本片；§6 起的正文属 s06a，本片不取。
% ============================================================

% ==== 原文 PDF 第 66 页（印刷页 66）====
\label{op:66}
% [说明: 本节首段的句子起于第 65 页末，第 66 页首接 "generation models…"；锚点置于该段之前。]

% [跨页接续: 原文第 65 页末→第 66 页首；"three | generation models" 跨页，s05b 止于上句 "This model is of interest for two reasons."]
首先，对簇 $X$ 取一个有限群的适当商，便给出已知最早、最简单的三代（three generation）模型之一，最初由 Schimmrigk [42] 构造。其次，Landau--Ginzburg 描述表明，该模型在其模空间的一个特定点上可精确求解。事实上，$S^{3}+ST^{3}$ 是 $E_{7}$ 模不变极小模型的超势，而 $S^{3}$ 是 $A_{1}$ 模不变极小模型的超势，因此该模型是三个 $E_{7}$ 极小模型与一个 $A_{1}$ 模型的张量积。它实际上就是 Gepner 的原始模型之一 [43]。

\subsection{拓扑的改变}

要点：本小节讨论：在弦论中，时空拓扑能否发生物理上连续的改变。答案是肯定的；局部机制是三维 Calabi--Yau 流形中的 flop（翻转），它对应一条有理曲线收缩、另一条有理曲线胀开的相变，而这一相变在 GLSM 相图中由 FI 参数 $r$ 的符号改变实现。

在弦论中，是否存在时空拓扑发生改变的物理过程，一直是一个长期令人着迷的问题。（借助对偶或镜像对称，人们已经知道一些不同拓扑之间的非经典等价性的例子，但那是略有不同的问题。）最近，Aspinwall、Greene 与 Morrison [7] 着手利用镜像对称来证明答案是肯定的，至少在一个特殊情形中如此。尽管他们的出发点相当不同，其分析最终涉及的正是本文一直在考察的那些相图。

一般而言，这些相图可能非常复杂，并且——如 [7] 在一个具体例子中详细展示的——可以包含多种 Calabi--Yau 相，对应不同的光滑靶空间。当 Fayet--Iliopoulos 项改变时，不同的靶时空之间会发生跃迁。按 §3.2 的论证，只要 $\theta$ 角取一般值，这些跃迁就是连续的。在弦论中，Fayet--Iliopoulos 参数（如同世界面拉格朗日量中的一切算符）被解释为时空中某些场的期望值；它们是动力学变量，可以随时间自由变化。这就给出了物理拓扑改变的一种机制。

以这种方式出现的跃迁，是拓扑不同但双有理等价的 Calabi--Yau 流形之间的跃迁。其理由在 §4 中已概述；简言之，通过改变 Fayet--Iliopoulos 参数对给定靶空间构造出的各种辛商，
% ==== 原文 PDF 第 67 页（印刷页 67）====
\label{op:67}
% [跨页接续: 原文第 66 页/第 67 页，句子跨页（原文在此断词 "iden-/tified"）]
一般地（在一个稠密开集上）都可以与同一靶空间的一个固定复商等同起来。

虽然 [7] 分析了一个整体情形，但作者这里的讨论范围要小一些，集中讨论拓扑改变的局部机制。换言之，我们将分析：通过改变 $D$ 函数中的常数项，如何在时空的某个局部区域得到拓扑的改变。读者应把我们将考虑的时空设想为嵌入在一个紧 Calabi--Yau 流形之中，而该流形的拓扑将在我们所考虑的过程中发生改变。

我们将考虑的例子处于复维数三，而且并不像看上去那样特殊。它就是 Mori 理论（复流形的双有理分类理论；综述见 [44]）在复维数三中考虑的、一般的“翻转”（flop，或曰不改变 $c_{1}$ 的双有理变换）。有可能通过引入足够多的参数，把三维 Calabi--Yau 流形之间的任何双有理变换构造成一连串变换，其中每一个都在局部同构于我们将分析的那一个。

我们考虑一个 $U(1)$ 规范场论，其中有两个电荷为 $1$ 的手征超场 $A_{i}$（$i=1,2$），以及两个电荷为 $-1$ 的手征超场 $B_{j}$（$j=1,2$）。可以把它们看作 $V\cong\mathbb{C}^{4}$ 上的坐标。取超势为零。玻色分量的通常势能是（在 $\sigma=0$ 时）
\begin{equation}
\tag{5.37}
U(a_{i},b_{j})=\frac{e^{2}}{2}\Biggl(\sum_{i}|a_{i}|^{2}-\sum_{j}|b_{j}|^{2}-r\Biggr)^{2}.
\end{equation}

$U(1)$ 在 $V$ 上的作用为
\begin{equation}
\tag{5.38}
\begin{aligned}
a_{i}&\to\lambda a_{i}\\
b_{j}&\to\lambda^{-1}b_{j}
\end{aligned}
\end{equation}
其中 $\lambda\in U(1)$；允许 $\lambda\in\mathbb{C}^{*}$ 便把它延拓为 $\mathbb{C}^{*}$ 作用。

% [Tutor Note: 推导补足。(5.37) 是解出辅助场 $D$ 之后的势能：由 §3 的 $D$ 方程，$D=-e^{2}\bigl(\sum_{i}|a_{i}|^{2}-\sum_{j}|b_{j}|^{2}-r\bigr)$，故 $U=\frac{1}{2e^{2}}D^{2}$。注意荷之和 $2\cdot1+2\cdot(-1)=0$：这既保证 $R$ 反常抵消（从而存在 $R$ 对称性），又保证 $\Theta$ 的 $\mathbb{C}^{*}$ 不变性（见后文）。$r$ 是 FI 参数，$r>0$、$r=0$、$r<0$ 三个区间给出不同的真空流形 $Z_{+},Z_{0},Z_{-}$；$\sigma$ 已被置零，它在 §3 的完整模型中可以非零。]

低能有效时空照例由令 $U=0$ 并对 $U(1)$ 取商得到；换言之，它就是 $V//U(1)$，即 $V$ 对
% ==== 原文 PDF 第 68 页（印刷页 68）====
\label{op:68}
% [跨页接续: 原文第 67 页/第 68 页，句子跨页；第 68 页自 "U(1) in the language of §4" 起]
$U(1)$ 的辛商——用 §4 的语言说。实际上有三种可能的辛商，分别对应 $r>0$、$r=0$ 与 $r<0$。我们称它们为 $Z_{+}$、$Z_{0}$ 与 $Z_{-}$。

按 §4 的论证，$Z_{+}$、$Z_{0}$ 与 $Z_{-}$ 在其自然复结构下都彼此双有理同构。事实上，令 $\tilde V$ 为 $V$ 中 $a_{i}$ 不全为零、且 $b_{j}$ 不全为零的区域。那么恰好按 §4 的论证，辛商 $\tilde V//U(1)$ 对任何 $r$ 都可以自然地与全纯商 $\tilde V/\mathbb{C}^{*}$ 等同起来。$Z_{+}$、$Z_{0}$ 与 $Z_{-}$ 是 $\tilde V/\mathbb{C}^{*}$ 的部分紧化，其做法是把某些坏的 $\mathbb{C}^{*}$ 轨道（即 $a_{1}=a_{2}=0$ 或 $b_{1}=b_{1}=0$ 的轨道）的贡献包含进来。
% [原文疑误: 此处原文印作 "b_1 = b_1 = 0"，据同句前半 "a_1 = a_2 = 0" 的对称形式、以及上文 "the b_j are not both zero" 的定义，应为 "b_1 = b_2 = 0"；建议改为 b_1 = b_2 = 0。]
事实上，这些坏轨道共有三种类型：原点 $O$（$a_{i}=b_{j}=0$）；由 $b_{j}=0$ 而 $a_{i}$ 不全为零的轨道构成的集合 $V_{1}$；以及由 $a_{i}=0$ 而 $b_{j}$ 不全为零的轨道构成的集合 $V_{2}$。

把适当的坏轨道包含进来，$Z_{+}$、$Z_{0}$ 与 $Z_{-}$ 便可描述如下：

(i) 对 $r>0$，不含 $D$ 的零点的 $\mathbb{C}^{*}$ 轨道是 $O$ 与 $V_{2}$ 中的轨道。其他 $\mathbb{C}^{*}$ 轨道（无论在 $\tilde V$ 中还是在 $V_{1}$ 中）各自恰好给辛商 $Z_{+}$ 贡献一个点。因此 $Z_{+}=(\tilde V\cup V_{1})/\mathbb{C}^{*}$。这里 $\tilde V\cup V_{1}$ 正是 $V$ 中 $a_{i}$ 不全为零的区域。$a_{i}$ 的取值在差一个 $\mathbb{C}^{*}$ 标度的意义下确定一个 $\mathbf{CP}^{1}$ 中的点，我们称它为 $\mathbf{CP}^{1}_{a}$。$Z_{+}$ 通过遗忘 $b_{j}$ 的取值而以 $\mathbf{CP}^{1}_{a}$ 为底做纤维化；由于 $b_{j}$ 的取值是任意的，纤维是 $\mathbb{C}^{2}$ 的一个拷贝。$Z_{+}\to\mathbf{CP}^{1}_{a}$ 的零截面，即轨迹 $b_{1}=b_{2}=0$，是一条亏格零的全纯曲线，实际上是一个嵌入 $\mathbf{CP}^{1}_{a}\subset Z_{+}$。此后我们把 $\mathbf{CP}^{1}_{a}$ 与它在此嵌入下的像等同起来。当 $b_{j}=0$ 时，$D$ 的消失给出 $\sum_{i}|a_{i}|^{2}=r$，因此 $\mathbf{CP}^{1}_{a}$ 的凯勒形式正比于 $r$。

(ii) 对 $r=0$，唯一有贡献的坏 $\mathbb{C}^{*}$ 轨道是 $O$；它给 $Z_{0}$ 贡献一个单点，而该点是一个奇点。我们稍后将更仔细地考察这个奇点。

(iii) $r<0$ 的结构与 $r>0$ 的相似，只是 $a_{i}$ 与 $b_{j}$ 的角色互换了。特别地，$Z_{-}=(\tilde V\cup V_{2})/\mathbb{C}^{*}$。这里 $\tilde V\cup V_{2}$ 是 $V$ 中 $b_{j}$ 不全为零的区域。$b_{j}$ 的取值确定
% ==== 原文 PDF 第 69 页（印刷页 69）====
\label{op:69}
% [跨页接续: 原文第 68 页/第 69 页，句子跨页]
$\mathbf{CP}^{1}$ 中的一个点，我们称它为 $\mathbf{CP}^{1}_{b}$。$Z_{-}$ 以 $\mathbf{CP}^{1}_{b}$ 为底做纤维化，纤维是 $\mathbb{C}^{2}$ 的一个拷贝，由 $a_{i}$ 的取值参数化。$Z_{-}\to\mathbf{CP}^{1}_{b}$ 的零截面 $a_{1}=a_{2}=0$ 给出 $Z_{-}$ 中的一条亏格零曲线，它是一个嵌入 $\mathbf{CP}^{1}_{b}\subset Z_{-}$。这条曲线——我们把它与 $\mathbf{CP}^{1}_{b}$ 等同——的凯勒形式正比于 $-r$。

% [Tutor Note: 物理图像。$r>0$ 时 $Z_{+}$ 是以 $\mathbf{CP}^{1}_{a}$ 为零截面、纤维为 $\mathbb{C}^{2}$ 的局部 Calabi--Yau 模型；$r<0$ 时 $\mathbf{CP}^{1}_{b}$ 取代之。$r\to0^{+}$ 时 $\mathbf{CP}^{1}_{a}$ 的面积正比于 $r$ 而趋于零，$r\to0^{-}$ 时 $\mathbf{CP}^{1}_{b}$ 的面积正比于 $-r$ 而趋于零。这正是 flop（翻转）的图像：一条有理曲线收缩为点，再沿另一个方向胀开；除去这一局部区域，复结构不变，因此 $c_{1}=0$ 得以保持。$\mathbb{C}^{*}$ 的“坏轨道”恰是使 $D=0$ 无解或解不唯一的那些轨道，把它们补进来就得到三个不同的部分紧化。]

$Z_{+}$、$Z_{0}$ 与 $Z_{-}$ 都是非紧的 Calabi--Yau 流形。事实上，$V$ 上的全纯四形式 $\Theta=\mathrm{d}a_{1}\wedge\mathrm{d}a_{2}\wedge\mathrm{d}b_{1}\wedge\mathrm{d}b_{2}$ 是 $\mathbb{C}^{*}$ 不变的（因为 $\sum_{i}Q_{i}=0$）；把它与生成 $\mathbb{C}^{*}$ 作用的向量场缩并，就得到一个处处非零的全纯三形式，它限制在 $V$ 的适当区域上正是 $Z_{+}$、$Z_{0}$ 或 $Z_{-}$ 上某个全纯体积形式的下拉。

% [Tutor Note: 推导补足。设 $\xi$ 为 $\mathbb{C}^{*}$ 作用的生成向量场，$\Omega=\iota_{\xi}\Theta$。$\Theta$ 是 $V\cong\mathbb{C}^{4}$ 上的体积形式，故与任何非零向量缩并都得到非零的 3-形式；又 $\Theta$ 的 $\mathbb{C}^{*}$ 不变性给出 $\mathcal{L}_{\xi}\Theta=0$，于是 $\mathcal{L}_{\xi}\Omega=0$ 且 $\iota_{\xi}\Omega=0$（水平），所以 $\Omega$ 可以下降为 $\tilde V/\mathbb{C}^{*}$（以及 $Z_{+},Z_{0},Z_{-}$ 的相应区域）上处处非零的全纯 3-形式——这正是 Calabi--Yau 条件（存在处处非零的全纯体积形式）的表述。注意 $\sum_{i}Q_{i}=0$ 正是 $\Theta$ 的 $\mathbb{C}^{*}$ 不变性所必需的。]

从 $r>0$ 过渡到 $r<0$ 时，拓扑发生了改变。当 $r$ 从上方趋于 $0$ 时，$\mathbf{CP}^{1}_{a}$ 收缩到零尺度；而当 $r<0$ 时，它被 $\mathbf{CP}^{1}_{b}$ 取代。经典地看，从 $r>0$ 到 $r<0$ 的插值经过 $r=0$ 处的一个奇点；但按 §3.2 的论证，只要 $\theta$ 角取一般值，$\sigma$ 模型中就没有这样的奇点。我们就在（非紧）Calabi--Yau 流形的 $\sigma$ 模型中实现了所许诺的光滑拓扑改变。

诚然，（通过交换 $a$ 与 $b$）$Z_{+}$ 与 $Z_{-}$ 实际上是同构的。然而，在整体情形中，当我们的讨论被应用于某个紧 Calabi--Yau 流形的一部分时，用 $Z_{-}$ 替换 $Z_{+}$、或用 $\mathbf{CP}^{1}_{b}$ 替换 $\mathbf{CP}^{1}_{a}$，就意味着时空拓扑的改变。

为了给该模型一个新的视角，现在更仔细地考察例外辛商 $Z_{0}$ 的原点处奇点的性质。设 $x$、$y$、$z$、$t$ 是 $\mathbb{C}^{4}$ 的一个拷贝上的坐标，我们称该拷贝为 $W$。诸公式
\begin{equation}
\tag{5.39}
\begin{aligned}
x&=a_{1}b_{1}\\
y&=a_{2}b_{2}\\
z&=a_{1}b_{2}\\
t&=a_{2}b_{1}.
\end{aligned}
\end{equation}

% ==== 原文 PDF 第 70 页（印刷页 70）====
\label{op:70}
% [跨页接续: 原文第 69 页/第 70 页，句子跨页（公式 (5.39) 之后接续）]
给出一个 $\mathbb{C}^{*}$ 不变的映射 $V\to W$。把这些公式限制到 $\tilde V\cup O$ 上再对 $\mathbb{C}^{*}$ 取商，便得到映射 $Z_{0}\to W$。（若从 $\tilde V\cup V_{1}$ 或 $\tilde V\cup V_{2}$ 出发，我们会得到映射 $Z_{\pm}\to W$。）显然，$Z_{0}$ 在 $W$ 中的像落在由
\begin{equation}
\tag{5.40}
xy-zt=0.
\end{equation}
定义的仿射二次超曲面 $Q$ 内。事实上，(5.39) 定义的映射是 $Z_{0}$ 与 $Q$ 之间的同构。

% [背景补充: $Q:\ xy-zt=0$ 就是三维复流形中著名的 conifold（锥形奇点）：它是 $\mathbb{C}^{4}$ 中的二次锥，在原点之外光滑。原文在这里不点名，但 (5.39) 中的 $x,y,z,t$ 正是 conifold 的四个坐标；两种“小消解”正是对它的两种小 blow-up（只胀开一条曲线、而不是一个除子）。]

该断言可如下论证：

(a) 为证明该映射是满射，必须证明：若 $x,y,z,t$ 满足 (5.40)，则 (5.39) 对 $a_{i},b_{j}$ 的某些取值成立。若 $x=y=z=t=0$，取 $a_{i}=b_{j}=0$。若比如 $x\neq0$，取 $a_{1}=1$，然后依次解 (5.39) 求出 $a_{2},b_{1},b_{2}$。

(b) 为证明该映射是单射，必须证明：若 $x,y,z,t$ 满足 (5.40)，则 (5.39) 在 $\mathbb{C}^{*}$ 作用的意义下唯一确定 $a_{i},b_{j}$。若 $x=y=z=t=0$，则 (5.39) 要求或者 $a_{i}$ 全为零、或者 $b_{j}$ 全为零。但在那种情况下，要使 $a_{i}$ 与 $b_{j}$ 定义 $Z_{0}$ 中的一个点，它们必须全为零，故 $x=y=z=t=0$ 只是 $Z_{0}$ 中点 $O$ 的像。否则，若比如 $x\neq0$，则 (5.39) 要求 $a_{1}\neq0$。我们可以用 $\mathbb{C}^{*}$ 作用设 $a_{1}=1$，然后 (5.39) 唯一确定 $a_{2},b_{1},b_{2}$。这就完成了 $Z_{0}$ 与 $Q$ 同构的证明。

$Q$ 在 $x=y=z=t=0$ 处的奇点是三维 Calabi--Yau 流形中最简单的一类孤立奇点（见 [41]，第 121 页）。把 $Q$ 的原点胀开可以消解这个奇点，但那样会破坏 Calabi--Yau 条件。在保持 Calabi--Yau 条件的前提下，有两种极小方式消解 $Q$ 的奇点。这两种选择是把奇点替换为 $\mathbf{CP}^{1}$ 的一个拷贝，而它应当就是上面构造的 $\mathbf{CP}^{1}_{a}$ 或 $\mathbf{CP}^{1}_{b}$。因此，$Q$ 作为 Calabi--Yau 流形的两种“小消解”（small resolutions）恰好就是 $Z_{+}$ 与 $Z_{-}$。我们上面找到的拓扑改变过渡，正是这两种小消解之间的过渡。

% [背景补充: 把 conifold 的原点直接胀开（blow-up）会得到一个例外除子，从而破坏 $c_{1}=0$；而“小消解”只用一条复曲线（这里是 $\mathbf{CP}^{1}$，其法丛为 $\mathcal{O}(-1)\oplus\mathcal{O}(-1)$）替换奇点，不引入除子，故保持 Calabi--Yau。两种小消解互为 flop：一种先把 $\mathbf{CP}^{1}_{a}$ 收缩到点，再沿另一方向胀出 $\mathbf{CP}^{1}_{b}$，反之亦然。]

% ==== 原文 PDF 第 71 页（印刷页 71）====
\label{op:71}

\subsubsection*{瞬子求和}

为了获得更多洞见，我们希望研究物理可观测量在从 $Z_{+}$ 到 $Z_{-}$ 的过渡中的行为；更确切地说，是从 $X_{+}$ 到 $X_{-}$ 的过渡，其中 $X_{\pm}$ 是紧 Calabi--Yau 流形，它们在一个区域（或可能有有限多个类似区域）之外彼此重合，而在该区域内它们看起来像 $Z_{\pm}$。最简单的可分析物理观测量是低能 Yukawa 耦合。由 $H^{2,1}(X_{\pm})$ 得到的无质量多重态的 Yukawa 耦合与 $r$ 无关，因此在从 $X_{+}$ 到 $X_{-}$ 的过渡下应当不变。更有意思的是与 $H^{1,1}(X_{\pm})$ 相联系的无质量多重态的 Yukawa 耦合。它们由复杂的瞬子求和决定；我们想知道，在从 $r>0$ 插值到 $r<0$ 时这些求和会变成什么样。

把 $X_{\pm}$ 上的完整瞬子求和写出来，需要对有理曲线有非常详细的信息。比较 $X_{\pm}$ 的瞬子则容易得多。在从 $X_{+}$ 到 $X_{-}$ 的过渡中，全部发生的事情就是：一条亏格零曲线——上面称为 $\mathbf{CP}^{1}_{a}$ 的那条——消失了，被另一条亏格零曲线——上面称为 $\mathbf{CP}^{1}_{b}$ 的那条——取代。我们的问题只是：计算 $r>0$ 时来自 $\mathbf{CP}^{1}_{a}$ 及其多重覆盖的贡献，并与 $r<0$ 时来自 $\mathbf{CP}^{1}_{b}$ 及其多重覆盖的贡献作比较。

首先，若 $H^{2,0}(X_{\pm})=0$（几乎所有的 Calabi--Yau 流形都如此），则 $H^{1,1}(X_{\pm})$ 可以等同于 $X_{\pm}$ 上的除子群（模线性等价）。再者，由于 $X_{+}$ 与 $X_{-}$ 只在复余维二上不同，$X_{+}$ 上的除子可以自然地与 $X_{-}$ 上的除子等同。这给出群 $H^{1,1}(X_{\pm})$ 之间的一个自然同构，因此我们干脆把它记作 $H^{1,1}(X)$。

看起来我们似乎需要关于 $H^{1,1}(X)$ 的某些详细信息，但幸运的是并非如此。若 $E_{1}$、$E_{2}$、$E_{3}$ 是三个除子，则有理曲线 $C$ 对相应 Yukawa 耦合的贡献正比于
\begin{equation}
\tag{5.41}
(C,E_{1})(C,E_{2})(C,E_{3}),
\end{equation}
其中 $(C,E)$ 是曲线 $C$ 与除子 $E$ 的交数。因此
% [Tutor Note: 这个公式来自世界面瞬子：绕 $C$ 的瞬子对三点耦合的贡献因子化到每条除子流上的顶点算符，给出三个交数的乘积；完整求和还需乘以瞬子权重 $e^{2\pi i\tau}$（见 (5.46) 前的讨论）与多重覆盖的因子。注意这里的 $E_{1},E_{2},E_{3}$ 是与 $H^{1,1}$ 中三个态对应的除子，$(C,E)$ 是流形 $X$ 上的拓扑交数，与 $r$ 无关。]

% ==== 原文 PDF 第 72 页（印刷页 72）====
\label{op:72}
% [跨页接续: 原文第 71 页/第 72 页，句子跨页]
$\mathbf{CP}^{1}_{a}$ 与 $\mathbf{CP}^{1}_{b}$ 及其覆盖只对与它们相交的那些除子所对应的态的 Yukawa 耦合有贡献。这意味着：任何除子，只要在线性等价意义下可以选取得不与 $\mathbf{CP}^{1}_{a}$ 和 $\mathbf{CP}^{1}_{b}$ 相交，就不受从 $X_{+}$ 到 $X_{-}$ 的插值影响。因此我们不必理解 $X_{\pm}$ 的除子类；只需理解 $Z_{\pm}$（等价地 $Z_{0}$）的除子类就够了。

$Z_{0}$ 的除子类群在 Hartshorne 的书 [45] 的例 II.6.6.1 与习题 II.6.5 中计算过。结果是：$Z_{0}$（或 $Z_{\pm}$）上的任何除子都是除子 $E$（由 $a_{1}=0$ 给出）的倍数。（$a_{1}$ 不是 $\mathbb{C}^{*}$ 不变的，因而是线丛的一个截面而不是函数，所以“$E$ 是 $a_{1}$ 的除子”这一事实并不使 $E$ 变成平凡除子。）例如，由 $b_{1}=0$ 给出的除子 $E'$ 满足
\begin{equation}
\tag{5.42}
E+E'=0
\end{equation}
因为乘积 $a_{1}b_{1}=x$ 是 $Z_{0}$ 上的一个函数，而 $E+E'$ 就是这个函数的除子。

% [Tutor Note: 推导补足。Witten 引用 Hartshorne 的例 II.6.6.1 与习题 II.6.5：$Z_{0}$（从而 $Z_{\pm}$）的除子类群由 $E=\{a_{1}=0\}$ 生成，即任何除子都是 $E$ 的倍数。微妙之处在于：$a_{1}$ 带 $U(1)$ 荷 $1$，在 $\mathbb{C}^{*}$ 商下不是函数而是某个线丛的截面，因此“$E$ 是 $a_{1}$ 的零点轨迹”并不使 $E$ 主化（$E$ 不是某个函数的除子）。另一方面，$x=a_{1}b_{1}$ 是 $\mathbb{C}^{*}$ 不变的多项式，在 $Z_{0}$ 上确为函数，故其除子 $E+E'$ 主化，即 $E'=-E$。于是 $\mathbf{CP}^{1}_{b}$ 对 $E$ 的交数可由它对 $E'$ 的交数反号得到——这正是 (5.44) 的来源。]

因此，我们唯一需要计算的 Yukawa 耦合，就是与 $E$ 对应的多重态的三点函数。为计算 $\mathbf{CP}^{1}_{a}$ 与 $\mathbf{CP}^{1}_{b}$ 对这个三点耦合的贡献，我们需要知道它们与 $E$ 的交数。的确，
\begin{equation}
\tag{5.43}
(\mathbf{CP}^{1}_{a},E)=1,
\end{equation}
因为 $\mathbf{CP}^{1}_{a}$ 与 $E$ 在一个点上横截相交（在 $V$ 上由 $\mathbb{C}^{*}$ 轨道 $a_{1}=b_{1}=b_{2}=0$、$a_{2}\neq0$ 代表）。同样地，$(\mathbf{CP}^{1}_{b},E')=1$，再考虑到 (5.42)，便推出
\begin{equation}
\tag{5.44}
(\mathbf{CP}^{1}_{b},E)=-1.
\end{equation}
（$E$ 与 $\mathbf{CP}^{1}_{b}$ 并不横截相交，因此这个数不能只靠数交点来确定。）

% [Tutor Note: 易错预警。(5.43) 可以逐点核对：$\mathbf{CP}^{1}_{a}$ 是 $Z_{+}$ 中 $b_{1}=b_{2}=0$ 的轨迹，而 $E$ 是 $a_{1}=0$；二者交于 $a_{1}=b_{1}=b_{2}=0$、$a_{2}\neq0$ 这一点，且相交横截。(5.44) 则不能靠数点：在 $Z_{-}$ 中 $\mathbf{CP}^{1}_{b}$ 是 $a_{1}=a_{2}=0$ 的轨迹，整条曲线都落在 $E=\{a_{1}=0\}$ 内，交不是真的横截交（曲线的交维数不为零），其交数只能通过 $E'=-E$ 与横截交数 $(\mathbf{CP}^{1}_{b},E')=1$ 间接得到。两条消解曲线对同一个除子的交数符号相反，这正是 flop 的拓扑指纹。]

% ==== 原文 PDF 第 73 页（印刷页 73）====
\label{op:73}

除了这些交数之外，为计算瞬子求和，还需要知道瞬子作用量：对一个瞬子 $C$，它贡献一个因子 $\exp(2\pi i\tau)$，其中 $-2\pi i\tau=\int_{C}\omega$；这里 $\omega$ 是代表 $X$ 上复化凯勒形式的二形式（它的实部是通常的凯勒形式，虚部包含 $\theta$ 角）。因此我们需要 $-2\pi i t_{a}=\int_{\mathbf{CP}^{1}_{a}}\omega$ 与 $-2\pi i t_{b}=\int_{\mathbf{CP}^{1}_{b}}\omega$。因为每个除子局部地都等价于 $E$ 的倍数，可以假定 $\omega$ 是 $E$ 的 Poincaré 对偶的倍数（再加上不贡献 $t_{a}$ 或 $t_{b}$ 的项）。考虑到 (5.43) 与 (5.44)，我们便得到重要的结果
\begin{equation}
\tag{5.45}
t_{b}=-t_{a}.
\end{equation}

% [Tutor Note: 推导补足。设 $\omega=m\,\mathrm{PD}(E)$（局部，略去不贡献 $t_{a},t_{b}$ 的项）。则 $-2\pi i t_{a}=\int_{\mathbf{CP}^{1}_{a}}\omega=m(\mathbf{CP}^{1}_{a},E)=m$，而 $-2\pi i t_{b}=\int_{\mathbf{CP}^{1}_{b}}\omega=m(\mathbf{CP}^{1}_{b},E)=-m$，故 $t_{b}=-t_{a}$。物理含义：flop 把该有理曲线所对应的那个复化凯勒模反号，于是 $X_{\pm}$ 的复化凯勒模空间的两支在 $t_{a}=0$ 处相接。]

现在我们要在 $X_{\pm}$ 上计算与除子 $E$ 相联系的三重场的 Yukawa 耦合。我们将把计算作为 $t_{a}$（或等价地 $t_{b}$）的函数来做。在 $X_{+}$ 上的计算中，物理区域是 $\mathrm{Im}\,t_{a}>0$，因为曲线 $\mathbf{CP}^{1}_{a}$ 若要存在（它在 $X_{+}$ 上确实存在），就必须具有正面积。$X_{-}$ 上计算的物理区域同样是 $\mathrm{Im}\,t_{b}>0$，等价地 $\mathrm{Im}\,t_{a}<0$。我们要做的是：对 $\mathrm{Im}\,t_{a}>0$ 计算 $X_{+}$ 上的瞬子求和，并把它解析延拓到 $\mathrm{Im}\,t_{a}<0$，与 $X_{-}$ 上的瞬子求和作比较。

实际计算并不困难。在 $X_{+}$ 上，我们必须计入 $\mathbf{CP}^{1}_{a}$ 及其多重覆盖。使用瞬子贡献的基本公式 [46] 以及由 Candelas 等人 [47] 猜想、并由 Aspinwall 与 Morrison [48] 论证的多重覆盖贡献公式，$\mathbf{CP}^{1}_{a}$ 及其多重覆盖对 Yukawa 耦合的贡献是
\begin{equation}
\tag{5.46}
\lambda_{+}=(\mathbf{CP}^{1}_{a},E)^{3}\sum_{n=1}^{\infty}e^{2\pi i n t_{a}}=\frac{e^{2\pi i t_{a}}}{1-e^{2\pi i t_{a}}}.
\end{equation}
第 $n$ 项是 $\mathbf{CP}^{1}_{a}$ 的第 $n$ 重覆盖的贡献；该级数在 $\mathrm{Im}\,t_{a}>0$ 时收敛。

% [Tutor Note: 推导补足。(5.46) 用了两件事：其一是第 $n$ 重覆盖对耦合的贡献带因子 $e^{2\pi i n t_{a}}$（$n$ 倍瞬子作用量）；其二是 Candelas 等人猜想、Aspinwall--Morrison 证明的多重覆盖公式，使求和成为几何级数 $\sum_{n\ge1}q^{n}=q/(1-q)$，其中 $q=e^{2\pi i t_{a}}$。收敛条件 $\mathrm{Im}\,t_{a}>0$ 即 $|q|<1$，对应 $\mathbf{CP}^{1}_{a}$ 有正面积。级数在 $t_{a}=0$ 处的极点位于两个物理区域的边界上，对应那条有理曲线面积趋于零的奇点。]

% ==== 原文 PDF 第 74 页（印刷页 74）====
\label{op:74}

在 $X_{-}$ 上，我们同样必须计入 $\mathbf{CP}^{1}_{b}$ 及其覆盖。与 (5.46) 类似的公式（此时在 $\mathrm{Im}\,t_{a}<0$ 时收敛）是
\begin{equation}
\tag{5.47}
\lambda_{-}=(\mathbf{CP}^{1}_{b},E)^{3}\sum_{n=1}^{\infty}e^{2\pi i n t_{b}}=-\frac{e^{-2\pi i t_{a}}}{1-e^{-2\pi i t_{a}}}.
\end{equation}

人们也许会天真地以为：从 $X_{+}$ 到 $X_{-}$ 光滑延拓的假设应当意味着 $\lambda_{+}=\lambda_{-}$；但事实上按上述公式
\begin{equation}
\tag{5.48}
\lambda_{+}-\lambda_{-}=-1.
\end{equation}

% [Tutor Note: 推导补足。代入 $u=e^{2\pi i t_{a}}$：$\lambda_{+}=u/(1-u)$；由 $(C_{b},E)^{3}=-1$ 与 $t_{b}=-t_{a}$，$\lambda_{-}=-u^{-1}/(1-u^{-1})=1/(1-u)$，故 $\lambda_{+}-\lambda_{-}=(u-1)/(1-u)=-1$，即 (5.48)。这两个表达式分别在 $\mathrm{Im}\,t_{a}>0$ 与 $\mathrm{Im}\,t_{a}<0$ 中收敛，互为穿过 $t_{a}=0$ 的解析延拓；常数差 $-1$ 不能由延拓本身消掉。]

不过，这一偏差有一个自然的解释。除了瞬子求和之外，Yukawa 耦合还从 $X_{+}$ 与 $X_{-}$ 的交配对（intersection pairing）得到常数式的“经典”贡献。双有理等价的流形 $X_{+}$ 与 $X_{-}$ 具有不同的上同调环和不同的交配对。必定是 $X_{+}$ 与 $X_{-}$ 的交配对之差为 $+1$，从而抵消 (5.48)。

这很容易验证。与其比较三重交数 $(E,E,E)$，我们同样可以考察 $(E,E',E')$，因为限制到 $V_{\pm}$ 上时 $E'=-E$。$X_{+}$ 与 $X_{-}$ 之间经典交配对的差，可以通过在 $V_{+}$ 与 $V_{-}$ 中数相交来度量。不过必须小心：在 $V_{+}$ 与 $V_{-}$ 中要对 $E,E',E'$ 使用相同的代表元；若把除子作线性等价的平移，则 $V_{\pm}$ 之外的相交数目可能会改变。于是我们用 $a_{1}=0$ 代表 $E$，用 $b_{1}=0$ 与 $b_{2}=0$ 代表两个 $E'$。在 $V_{+}$ 中，这些除子横截相交于唯一一点 $a_{1}=b_{1}=b_{2}=0$，故交数为 $+1$。在 $V_{-}$ 中，$b_{1}$ 与 $b_{2}$ 从不同时为零，故交数为 $0$。于是我们得到预期的、$X_{+}$ 相对于 $X_{-}$ 的超出交数（excess intersection number）$+1$。

% [Tutor Note: 物理图像。瞬子部分差 $-1$、经典交数部分差 $+1$，两者相加使总 Yukawa 耦合在 flop 下连续——这正是弦论拓扑改变“光滑”的具体体现。另一方面，$X_{+}$ 与 $X_{-}$ 的上同调环不同（三重交数相差 $1$），说明它们确实是拓扑不同的流形；$B$ 模型（复结构一侧）在 flop 下不变，$A$ 模型（凯勒一侧）则通过 $t_{a}\to-t_{a}$ 的解析延拓相联系。]

所以，$X_{-}$ 上的 Yukawa 耦合正如所预期的那样，是 $X_{+}$ 上 Yukawa 耦合的解析延拓。此外，(5.46) 与 (5.47) 在 $t_{a}=0$ 处有一个极点，且这是它们唯一的奇点；这一事实
% ==== 原文 PDF 第 75 页（印刷页 75）====
\label{op:75}
% [跨页接续: 原文第 74 页/第 75 页，句子跨页；本片仅收 §5.5 末句 3 行，其后即 "6." 标题，属 s06a]
与 §3.2 的论证相符：按该论证，$r$ 与 $\theta$（本质上分别是 $t_{a}$ 的虚部与实部）必须同时为零，才会出现奇点。

% [Tutor Note: 物理图像。$t_{a}=0$ 即 $r=0$ 且 $\theta=0$，这是两个相的分界点，也是唯一可能出现奇点之处；$\theta$ 取一般值时，$r$ 穿越 $0$ 的过渡在 $\sigma$ 模型中是光滑的。这正是“拓扑改变”作为物理过程的含义：时空几何的拓扑在 $r$ 的符号改变时改变，但弦论本身没有奇异行为——奇点只出现在经典几何的描述中（$Z_{0}$ 的 conifold 奇点）。]
